\section{The Mutable-\cpi{} Baseline}\label{sec:baseline}

In this section we describe \cbnd{}~\cite{NuclearCD}, the state-of-the-art exact algorithm for general $(r,s)$-nucleus decomposition.
%
We first walk through its peeling loop, then enumerate the residual structures it maintains, and finally pinpoint the operations that this paper specialises away when restricted to $r=1$.

\stitle{The peeling loop.}
\cbnd{} follows the standard core-decomposition template~\cite{coreAlg}: it maintains the residual support of every alive $r$-clique in a priority queue, repeatedly pops the $r$-clique of minimum residual support, assigns it the current peel level as its core value, and propagates the support loss to the remaining alive $r$-cliques.
%
At $r=1$ the queue items are vertices, so each iteration pops a vertex $v$, records $\kappa_s(v)$, and decrements the residual support of every vertex that loses an $s$-clique witness because of $v$'s removal.
%
The non-trivial work happens during this support-loss propagation, because the \cpi{} that encodes the witnesses must be brought into a state consistent with $v$ being gone before the next pop.

\stitle{Residual state.}
To keep the witness counts consistent across peels, \cbnd{} maintains the \cpi{} as mutable residual state.
%
For each path $\TPath$ it stores the hold and pivot vertex sets as hash sets so that the test ``does $v$ belong to this path?'' is constant time.
%
A separate reverse map sends every vertex $u$ to the list of paths that contain it, so that when $u$ is popped the algorithm can find the affected paths without scanning all of $\mathcal{P}$.
%
A heap (or bucket queue) keeps every alive vertex ordered by current residual support, with decrease-key operations for the support drops triggered at each peel event.
%
Finally, the Bron--Kerbosch recursion tree that produced the \cpi{} is kept mutable: when a peel event invalidates a path, the algorithm walks into the tree and rewrites or removes the affected nodes.

\stitle{Per-peel work.}
The following proposition records the cost of one peel event under the residual layout above; it is a direct consequence of the pseudocode in~\cite{NuclearCD} restricted to $r=1$, and is the quantity that the static-\cpi{} invariant of Section~\ref{sec:static} replaces by counter arithmetic.

\begin{proposition}[\cbnd{} Per-Peel Cost at $r=1$]\label{prop:cbnd-cost}
Let \cbnd{} peel a vertex $v$ at residual support $k$. The event performs, for each path $\TPath$ containing $v$:
(i) an $O(1)$ hash query to decide whether $v\in\Vh(\TPath)$ or $v\in\Vp(\TPath)$;
(ii) a path rewrite that erases the reverse-map entries of every other vertex in $\Vh(\TPath)\cup\Vp(\TPath)$, in time $\Theta(|\Vh(\TPath)|+|\Vp(\TPath)|)$;
(iii) a support recompute over the surviving encoded $s$-cliques on $\TPath$ followed by a heap decrease-key for every affected path-mate.
%
Summed over paths through $v$, the per-peel cost is $\Theta\bigl(M_{\TPath}(v)+R(v)+\log n\bigr)$, where $M_{\TPath}(v)=\sum_{\TPath\ni v}\!|\Vh(\TPath)\cup\Vp(\TPath)|$ is the vertex--path incidence touched and $R(v)$ is the total support-recompute cost for $v$'s path-mates.
%
The whole peel therefore runs in $\Theta(\Sig+R+n\log n)$ time, with $R=\sum_v R(v)$.
\end{proposition}

\begin{proof}
Items~(i)--(iii) read off the pseudocode of~\cite{NuclearCD} specialised to $r=1$: the hash query is a single reverse-map lookup; rewriting $\TPath$ visits each $u\in\Vh(\TPath)\cup\Vp(\TPath)\setminus\{v\}$ once to erase the $(u,\TPath)$ reverse-map entry, so the rewrite is $\Theta(|\Vh(\TPath)\cup\Vp(\TPath)|)$; the recompute and heap decrease-key together cost $R$ summed over the peel plus one $O(\log n)$ priority-queue operation per popped vertex.
%
Summing (i)--(iii) over the $\deg_{\Sig}(v)$ paths through $v$ gives the per-peel bound, and aggregating over all vertices gives the global bound by $\sum_v M_{\TPath}(v)=\Sig$.
%
\qed
\end{proof}

The two dominant components are $M_{\TPath}$ (vertex--path incidence rewrites) and $R$ (support recomputes); both are intrinsic to maintaining the \cpi{} as mutable state.

\stitle{Why this state is needed in the general case.}
The mutable design is not over-engineering for the general $(r,s)$ setting.
%
When the peeled object is an edge ($r=2$) or a higher-order clique ($r\ge3$), the removal can destroy edges among vertices that remain alive; some of the cliques that a path encodes therefore cease to exist as cliques.
%
In that setting, the stored hold and pivot sets of $\TPath$ no longer define a valid encoded clique family, so the path must be split, shrunk, or deleted by walking back into the recursion tree.
%
The reverse map and the per-path hash sets exist precisely to support that surgery efficiently.

\stitle{Why this state is wasteful at $r=1$.}
The $(1,s)$ case has a stronger invariant: vertex peeling removes vertices but does not delete edges, so a set of vertices that formed a clique when the \cpi{} was built remains a clique throughout the peel (we prove this in Section~\ref{sec:static}).
%
A path's encoded clique family therefore stays structurally valid; only the question of whether a particular vertex of $\Vh(\TPath)\cup\Vp(\TPath)$ is still alive changes.
%
This means the path rewrites, the reverse-map erases, and the recursion-tree surgery in the per-peel cost above are all unnecessary at $r=1$ — they maintain a mutable residual state that is logically equivalent to a static incidence layout plus a single integer counter per path.
%
The next four sections cash in this observation: Section~\ref{sec:static} states the static-\cpi{} invariant in counter form; Section~\ref{sec:hindex} derives \spin{} and \spinStar{} over the resulting static index; Section~\ref{sec:parabuild} parallelises construction once peeling becomes counter-only; and Section~\ref{sec:buildhier} extracts the nucleus hierarchy from the preserved \cpi{} without any clique re-enumeration.
